% =========================================
% CHAPTER 7: SECURITY ANALYSIS
% =========================================
\chapter{Security Analysis and Failure Modes}
\label{ch:security}

This chapter examines security properties and failure modes. Two of the scenarios below were directly tested: sequencer crashes (scenario 3) and invalid state submissions (scenario 1). The remaining three---chain reorganisations, censorship, and data unavailability---are not easily reproducible on a local Hardhat node, so they are analysed from the architecture rather than measured. Each section makes clear which category it falls into.

\section{Threat Model}

Five adversarial scenarios are relevant to any Layer~2 deployment:

\begin{enumerate}
    \item \textbf{Malicious sequencer}: the sequencer attempts to commit invalid state transitions to L1.
    \item \textbf{L1 chain reorganization}: one or more confirmed L1 blocks are replaced, invalidating L2 commitments anchored to those blocks.
    \item \textbf{Sequencer failure}: the sequencer crashes or becomes temporarily unavailable.
    \item \textbf{Censorship attack}: the sequencer selectively excludes certain transactions from batches.
    \item \textbf{Data unavailability}: L2 state data is not published or becomes inaccessible, preventing independent verification.
\end{enumerate}

We assume the adversary controls the sequencer process but not the Layer~1 consensus mechanism. The L1 chain is assumed to be secure under standard proof-of-stake assumptions---that is, at least two-thirds of validators are honest. Each section below examines how each architecture responds to these threats, with reference to the empirical results from Chapter~\ref{ch:evaluation} where applicable.

\section{ZK Rollup Security}

\subsection{Strength: Cryptographic Validity}

The security of a ZK~Rollup derives from the mathematical properties of the underlying proof system:
\begin{itemize}
    \item \textbf{Soundness}: a valid proof guarantees that the corresponding state transition is correct; it is computationally infeasible to forge a proof for an invalid transition.
    \item \textbf{Zero-knowledge}: the proof reveals no information about the individual transactions within the batch, only that the transition is valid.
    \item \textbf{Non-interactivity}: proof verification is a single-step operation that requires no further communication between the prover and the verifier.
\end{itemize}

The practical implication is that any invalid state is rejected at the point of L1 verification. Invalid commitments never enter L1 storage, providing a strong safety guarantee.

\subsection{Vulnerability: Prover Availability}

If the sequencer produces valid transactions but fails to generate proofs (whether due to a crash, resource exhaustion, or deliberate refusal), the system enters a degraded state:
\begin{itemize}
    \item transactions continue to be accepted at the L2 level but are not finalized on L1;
    \item users cannot prove the current L2 state to external parties or to L1;
    \item the system is effectively stalled but not compromised---no funds are at risk, but progress halts.
\end{itemize}

The standard mitigation is a force-inclusion mechanism that requires the sequencer to produce proofs within a deadline or face penalties (such as collateral slashing or sequencer replacement).

\subsection{Vulnerability: Proof System Compromise}

If the zero-knowledge proof system itself contains a mathematical flaw, it may become possible to construct proofs for invalid state transitions. This is a low-probability but high-impact risk. Standard mitigations include:
\begin{itemize}
    \item adopting well-studied, battle-tested proof systems such as PLONK or STARK;
    \item commissioning regular security audits by cryptography experts;
    \item monitoring the academic literature for newly discovered vulnerabilities in the underlying assumptions.
\end{itemize}

\section{Optimistic Rollup Security}

\subsection{Strength: Economic Incentives}

The security of an Optimistic~Rollup is grounded in economic game theory rather than cryptographic proofs:
\begin{itemize}
    \item \textbf{Collateral requirements}: both the sequencer and any challenger must post bonds before participating in the protocol.
    \item \textbf{Slashing}: a party found to be dishonest (either a malicious sequencer or a frivolous challenger) loses its posted collateral.
    \item \textbf{Rational incentives}: under standard game-theoretic assumptions, honest behavior is the economically dominant strategy.
\end{itemize}

This design allows invalid states to be challenged and reverted, but only through an economic dispute process rather than an automatic cryptographic check.

\subsection{Vulnerability: Challenger Absence}

The most significant vulnerability of an Optimistic~Rollup is the possibility that no honest challenger submits a fraud proof during the challenge period. This can occur under the following conditions:
\begin{itemize}
    \item low network participation, with few nodes running full archival copies of the L2 state;
    \item prohibitively high challenge costs that exceed the economic incentive to challenge;
    \item coordinated collusion between the sequencer and a majority of potential challengers.
\end{itemize}

We can model the probability of a successful attack (i.e., no challenge) as follows. Let~$p$ denote the probability that any single observer detects the fraud, and let~$n$ denote the number of independent observers. Then:
\begin{equation}
P(\text{no challenge}) = (1 - p)^n
\end{equation}

For example, with $p = 0.1$ and $n = 10$:
\begin{equation}
P(\text{no challenge}) = (0.9)^{10} \approx 0.35
\end{equation}

This illustrative calculation shows that even with 10 independent observers, each having a 10\% chance of detecting fraud, there remains a 35\% probability that no challenge is submitted. In practice, production systems mitigate this through dedicated watchtower services and by increasing the number of independent validators.

More generally, we can express the minimum number of observers required to reduce the attack probability below a threshold~$\epsilon$ as:
\begin{equation}
n \geq \frac{\ln \epsilon}{\ln (1 - p)}
\end{equation}

For $p = 0.1$ and $\epsilon = 0.01$, we obtain $n \geq \frac{\ln 0.01}{\ln 0.9} \approx 43.7$, meaning at least 44~independent observers are needed to reduce the probability of a successful unchallenged attack to below 1\%. This requirement places a concrete lower bound on the validator set size for meaningful Optimistic security.

\subsection{Vulnerability: Withdrawal Latency}

Users who wish to withdraw funds from an Optimistic~Rollup to L1 must wait for the full challenge period (7~days in production) to elapse. This has several practical consequences:
\begin{itemize}
    \item users who need faster access to their funds must rely on third-party liquidity providers, who charge a fee for the service;
    \item withdrawals of small amounts become disproportionately expensive relative to liquidity provider fees;
    \item the user experience is significantly worse than that of a ZK~Rollup, where withdrawals can be processed as soon as the proof is verified on L1.
\end{itemize}

Proposed mitigations include liquidity networks and hybrid models that use ZK proofs to accelerate the withdrawal process.

\section{L1 Reorganization Impact}

The following analysis is derived from the architecture, not from an experiment. Hardhat produces deterministic blocks and does not simulate chain reorganisations, so these scenarios were not run. The estimates use Ethereum mainnet block timing (12 seconds per block) as a reference.

\subsection{Scenario: Shallow Reorg (1~Block)}

Consider the following situation: an L1 block at height~$n$ contains a ZK proof~$\pi_1$ for batch~$B_1$. A 1-block reorganization occurs, and block~$n$ is replaced by a new block that does not contain~$\pi_1$.

The impact on the ZK~Rollup is as follows:
\begin{itemize}
    \item the proof $\pi_1$ is no longer part of the L1 canonical chain;
    \item the L2 state update that was based on $\pi_1$ becomes invalid;
    \item the system reverts to the previous valid state;
    \item the sequencer must resubmit the batch to L1.
\end{itemize}

The recovery process is automatic:
\begin{enumerate}
    \item The sequencer detects the reorganization by monitoring the L1 chain.
    \item It regenerates the proof (possibly with updated batch ordering if necessary).
    \item It resubmits the proof to L1.
    \item The system re-finalizes once the new proof is included and confirmed.
\end{enumerate}

The expected recovery time can be estimated as:
\begin{equation}
T_{\text{recovery}} = T_{\text{detection}} + T_{\text{reproof}} + T_{\text{re-inclusion}}
\end{equation}

For a 1-block reorg on Ethereum: detection~$\approx$~12s, reproof~$\approx$~95ms, re-inclusion~$\approx$~12s, yielding a total recovery time of approximately \textbf{24~seconds}. Importantly, no data is lost during this process.

\subsection{Scenario: Deep Reorg ($>$7~Blocks)}

A deep reorganization ($>$7~blocks) is an extremely unlikely event on Ethereum, particularly following the transition to proof-of-stake and the introduction of finality gadgets. Such an event would require either a successful 51\% attack or a catastrophic network partition.

If a deep reorg were to occur:
\begin{itemize}
    \item both ZK and Optimistic~Rollups would need to resubmit multiple batches or state roots;
    \item in the Optimistic case, economic incentives could break down if challengers are unable to participate during the disruption;
    \item the system would likely require manual intervention or social consensus to recover.
\end{itemize}

For practical purposes, deep reorgs are not considered a realistic threat for established L1 networks.

\section{Comparative Security Summary}

Table~\ref{tab:security-summary} provides a qualitative summary of the security properties of each architecture across the dimensions considered in this chapter.

\begin{table}[h]
    \centering
    \begin{tabular}{l|c|c}
        \hline
        \textbf{Security Property} & \textbf{ZK Rollup} & \textbf{Optimistic} \\
        \hline
        Cryptographic safety & Strong & Economic \\
        \hline
        Recovery speed & Moderate & Moderate \\
        \hline
        Censorship resistance & Good & Medium \\
        \hline
        Liveness requirement & Prover only & Prover + challenger \\
        \hline
        Finality speed & Fast ($\sim$15s) & Slow (7~days) \\
        \hline
        User exit cost & Low & High (7~days or fee) \\
        \hline
    \end{tabular}
    \caption{Comparative security summary}
    \label{tab:security-summary}
\end{table}

\section{Data Availability Analysis}

Data availability---the guarantee that all L2 transaction data is accessible to any observer---is a prerequisite for the security of both architectures, though for different reasons.

\subsection{Why Data Availability Matters}

For a \textbf{ZK~Rollup}, data availability enables independent state reconstruction. Even though the validity proof guarantees that the state transition is correct, users need access to the underlying transaction data to determine their own account balances, construct withdrawal proofs, and verify that their transactions were included. Without data availability, the ZK~Rollup reduces to a system where correctness is assured but accessibility is not---users must trust the sequencer to report their state truthfully.

For an \textbf{Optimistic~Rollup}, data availability is even more critical. Challengers cannot construct a fraud proof without access to the transaction data from the disputed batch. If the sequencer withholds data, it becomes impossible to prove that a state root is incorrect, effectively neutralizing the entire challenge mechanism. A data-withholding sequencer in an Optimistic~Rollup can commit arbitrary state transitions with impunity.

\subsection{Data Availability in Our Framework}

Both sequencers in our framework post full batch data as L1 calldata, ensuring that any Ethereum full node can reconstruct the complete L2 state from L1 data alone. The cost of this approach is captured in the gas analysis of Chapter~\ref{ch:evaluation}: each byte of calldata costs 16~gas (or 4~gas for zero bytes), which constitutes a significant fraction of the per-submission gas cost.

Alternative data availability strategies---off-chain data with on-chain commitments (Validium), or per-transaction data availability choices (Volition)---offer lower costs but weaken the data availability guarantee. Our framework does not implement these alternatives, but they represent a natural direction for extension (Chapter~\ref{ch:conclusion}).

\section{Practical Implications}

The security analysis yields several practical guidelines for deploying and operating Layer~2 systems.

\textbf{Validator set requirements.} The challenger absence analysis demonstrates that Optimistic~Rollups require a minimum validator set size to achieve a target security level. Operators should monitor the number of active watchtowers and ensure diversity across jurisdictions and infrastructure providers to prevent correlated failures.

\textbf{Reorg monitoring.} Both architectures benefit from active monitoring of L1 chain reorganizations. The ZK~Rollup's automatic recovery (24~seconds for a 1-block reorg) is faster than any manual intervention, but operators should still alert on reorg events to track their frequency and depth. Optimistic~Rollups are more vulnerable during deep reorgs, as the challenge mechanism may be disrupted.

\textbf{Sequencer redundancy.} The crash recovery results from Chapter~\ref{ch:evaluation} show that both sequencers lose in-memory state upon failure. Production deployments should implement transaction persistence (write-ahead logging) and standby sequencer failover to minimize data loss during outages.

\textbf{Cost of security.} The ZK~Rollup's security comes at a measurable cost: approximately 2.6$\times$ higher L1 gas consumption in the Hardhat test environment (Chapter~\ref{ch:evaluation}). For applications where the security premium is justified by the value at risk, this cost is acceptable. For lower-value applications, the Optimistic model's cheaper security may be sufficient.

\section{What the Security Analysis Adds}

The patterns from Chapter~\ref{ch:evaluation} translate directly into security terms. ZK finality is fast because invalid states are rejected at the contract level before they reach L1 storage. Optimistic finality is slow because security depends on watchers having enough time to act. The analysis in this chapter gives that intuition a concrete form---specifically, the minimum observer count needed to bring unchallenged-fraud probability below a target threshold, and the recovery-time formula for reorg scenarios.

Chapter~\ref{ch:conclusion} takes these findings and asks the practical question: when does each architecture actually make sense, and what would a follow-on study need to test next?
